\section{From residual estimation to execution recovery}
\label{sec:theory}\label{sec:contraction}
We analyze the evaluated prediction-only updates in the order in which
errors reach the robot: selected-coordinate estimation, the applied
correction, a scheduled execution tube, and predictor--execution coupling.
The nominal comparison follows the \emph{same supplied VLA commands}.
Its contraction metric concerns the robot with its low-level controller;
neither the VLA Jacobian nor a separately evolving healthy policy is required.
The metric may be constructed from physical tracking dynamics
(Appendix~\ref{app:gen-innerloop}) and, for the evaluated robots, was
measured by same-command replay; that measurement is what decides whether
the results below are certificates or finite-horizon statements
(Appendix~\ref{app:measured-certificates}). Prediction-only adaptation does
not compute the nominal response or a geodesic online.

The execution state must close the physical transition: it contains robot
and necessary low-level controller, delay, and interpolation states, plus
object/contact states when their evolution affects execution. Omitted contact
or environment effects require a uniform disturbance bound on the stated
domain; otherwise the certificate applies only to the reduced servo model
(Appendix~\ref{app:physical-closure}). VLA internal memory and the separate
observer FIR history are excluded unless a retained physical dependence
requires additional state or bounded forcing
(Appendix~\ref{app:history-partition}).

Same-input contraction and disturbance-to-tracking bounds are established
machinery \citep{lohmiller1998,manchester2017}. Our
specialization composes residual calibration with partial correction support,
update holds, snapshot age, clipping and realization error, and a
prediction-dependent small-gain condition. The optional continuous composite
controller and fixed-data gradient candidate are developed only in
Appendices~\ref{sec:composite-extension} and~\ref{app:sampled-gradient}.
